
\documentclass[a4paper,12pt]{article}
\usepackage[italian]{babel}
\usepackage{geometry}
\usepackage{listings}
\geometry{margin=2.5cm}

\title{Soluzioni - Verifica SQL Sistema Universitario}
\date{}

\begin{document}
\maketitle

\section*{Soluzioni Query}

\begin{enumerate}
	
	\item
	\begin{verbatim}
		SELECT nome, cognome
		FROM STUDENTE
		ORDER BY cognome, nome;
	\end{verbatim}
	
	\item
	\begin{verbatim}
		SELECT C.nomeCorso, C.crediti, D.cognome
		FROM CORSO C
		JOIN DOCENTE D ON C.idDocente = D.idDocente
		ORDER BY C.crediti DESC;
	\end{verbatim}
	
	\item
	\begin{verbatim}
		SELECT DISTINCT D.nome, D.cognome
		FROM DOCENTE D
		JOIN CORSO C ON D.idDocente = C.idDocente
		WHERE D.dipartimento = 'Informatica';
	\end{verbatim}
	
	\item
	\begin{verbatim}
		SELECT nome, cognome, annoIscrizione
		FROM STUDENTE
		WHERE citta = 'Roma'
		ORDER BY annoIscrizione;
	\end{verbatim}
	
	\item
	\begin{verbatim}
		SELECT C.nomeCorso, C.crediti, D.cognome
		FROM CORSO C
		JOIN DOCENTE D ON C.idDocente = D.idDocente
		WHERE C.crediti > 6;
	\end{verbatim}
	
	\item
	\begin{verbatim}
		SELECT S.nome, S.cognome, C.nomeCorso
		FROM ISCRIZIONE I
		JOIN STUDENTE S ON I.idStudente = S.idStudente
		JOIN CORSO C ON I.idCorso = C.idCorso;
	\end{verbatim}
	
	\item
	\begin{verbatim}
		SELECT S.nome, C.nomeCorso, I.voto
		FROM ISCRIZIONE I
		JOIN STUDENTE S ON I.idStudente = S.idStudente
		JOIN CORSO C ON I.idCorso = C.idCorso;
	\end{verbatim}
	
	\item
	\begin{verbatim}
		SELECT C.nomeCorso, COUNT(*) AS numero_studenti
		FROM ISCRIZIONE I
		JOIN CORSO C ON I.idCorso = C.idCorso
		GROUP BY C.idCorso, C.nomeCorso;
	\end{verbatim}
	
	\item
	\begin{verbatim}
		SELECT C.nomeCorso, AVG(I.voto) AS voto_medio
		FROM ISCRIZIONE I
		JOIN CORSO C ON I.idCorso = C.idCorso
		GROUP BY C.idCorso, C.nomeCorso;
	\end{verbatim}
	
	\item
	\begin{verbatim}
		SELECT S.nome, S.cognome
		FROM STUDENTE S
		LEFT JOIN ISCRIZIONE I
		ON S.idStudente = I.idStudente
		WHERE I.idIscrizione IS NULL;
	\end{verbatim}
	
	\item
	\begin{verbatim}
		SELECT C.nomeCorso
		FROM CORSO C
		LEFT JOIN ISCRIZIONE I
		ON C.idCorso = I.idCorso
		WHERE I.idIscrizione IS NULL;
	\end{verbatim}
	
	\item
	\begin{verbatim}
		SELECT D.cognome, COUNT(*) AS numero_corsi
		FROM DOCENTE D
		JOIN CORSO C ON D.idDocente = C.idDocente
		GROUP BY D.idDocente, D.cognome
		HAVING COUNT(*) > 1;
	\end{verbatim}
	
	\item
	\begin{verbatim}
		SELECT S.nome, S.cognome, COUNT(*) AS numero_esami
		FROM STUDENTE S
		JOIN ISCRIZIONE I ON S.idStudente = I.idStudente
		GROUP BY S.idStudente, S.nome, S.cognome
		HAVING COUNT(*) >= 2;
	\end{verbatim}
	
	\item
	\begin{verbatim}
		SELECT C.nomeCorso, AVG(I.voto) AS voto_medio
		FROM ISCRIZIONE I
		JOIN CORSO C ON I.idCorso = C.idCorso
		GROUP BY C.idCorso, C.nomeCorso
		HAVING AVG(I.voto) > 25;
	\end{verbatim}
	
	\item
	\begin{verbatim}
		SELECT idStudente, COUNT(*) AS numero_esami
		FROM ISCRIZIONE
		GROUP BY idStudente
		ORDER BY numero_esami DESC
		LIMIT 1;
	\end{verbatim}
	
	\item
	\begin{verbatim}
		SELECT D.cognome, COUNT(*) AS numero_studenti
		FROM DOCENTE D
		JOIN CORSO C ON D.idDocente = C.idDocente
		JOIN ISCRIZIONE I ON C.idCorso = I.idCorso
		GROUP BY D.idDocente, D.cognome
		HAVING COUNT(*) >= 3;
	\end{verbatim}
	
	\item
	\begin{verbatim}
		SELECT C.nomeCorso
		FROM CORSO C
		LEFT JOIN ISCRIZIONE I
		ON C.idCorso = I.idCorso
		WHERE I.dataEsame IS NULL;
	\end{verbatim}
	
	\item
	\begin{verbatim}
		SELECT C.nomeCorso, AVG(I.voto) AS voto_medio
		FROM ISCRIZIONE I
		JOIN CORSO C ON I.idCorso = C.idCorso
		GROUP BY C.idCorso, C.nomeCorso
		ORDER BY voto_medio DESC
		LIMIT 1;
	\end{verbatim}
	
	\item
	\begin{verbatim}
		SELECT DISTINCT S.nome, S.cognome
		FROM STUDENTE S
		JOIN ISCRIZIONE I ON S.idStudente = I.idStudente
		WHERE I.voto = 30;
	\end{verbatim}
	
	\item
	\begin{verbatim}
		SELECT S.nome, S.cognome
		FROM STUDENTE S
		JOIN ISCRIZIONE I ON S.idStudente = I.idStudente
		GROUP BY S.idStudente, S.nome, S.cognome
		HAVING COUNT(DISTINCT I.idCorso) >= 2;
	\end{verbatim}
	
\end{enumerate}

\end{document}
